% Fichier complété par tools/complete_missing_corrections.py.
% Source : CIN/CIN-02-VitesseAcceleration/06_TR/06_TR.tex
% Statut : rédigé (4 question(s)).
\begin{corrigebox}[Corrigé rédigé]
\CorrigeQuestion{1}
Le point $B$ est l'origine du pivot 2/1 et appartient au coulisseau 1 : il reste donc sur la droite $(A,\vec i_0)$. L'ensemble accessible est le segment ou la droite correspondant au domaine de $\lambda$.
\CorrigeQuestion{2}
$\overrightarrow{AB}=\lambda(t)\vec i_0$, donc $x_B=\lambda(t)$, $y_B=z_B=0$. Si le point visé est l'extrémité $C$, alors $\overrightarrow{AC}=\lambda\vec i_0+R(\cos\theta\vec i_0+\sin\theta\vec j_0)$.
\CorrigeQuestion{3}
Pour imposer une trajectoire $C(t)=(x(t),y(t))$, on résout
                $x=\lambda+R\cos\theta$, $y=R\sin\theta$. Ainsi
                $\theta(t)=\arcsin(y(t)/R)$ (branche choisie selon la configuration) et
                $\lambda(t)=x(t)-R\cos\theta(t)$. Pour un segment parcouru à vitesse $v$,
                $x(t)=x_0+vt$ et $y(t)=y_0$.
\CorrigeQuestion{4}
Un tracé Python direct utilise un tableau de temps, calcule
                `theta=np.arcsin(y/R)` puis `lam=x-R*np.cos(theta)`, et vérifie la
                trajectoire avec `xc=lam+R*np.cos(theta)` et `yc=R*np.sin(theta)`.
\end{corrigebox}
